\section{Development Timeline}
\label{sec:timeline}

The repository's git history spans \textbf{27 May 2026} (initial commit) to
\textbf{22 July 2026} (most recent committed work at the time of this
report), across 57 commits organized into five working sessions of
increasing scope. This section gives the narrative account of each session;
Section~\ref{sec:tech-decisions} groups the underlying bug fixes and design
decisions by theme rather than chronology, for readers who want the ``why''
without the session structure.

\subsection{Pre-history: bring-up and the VLA-to-LLM/VLM pivot (late May 2026)}

The earliest commits (\texttt{5478d4e}, \texttt{245b34d}, \texttt{3484aa9})
establish the ROS2 package skeleton, including refactoring an early
monolithic \texttt{robo\_reason\_real} package into the seven focused
packages listed in Section~\ref{sec:architecture}, and rewriting RG2 gripper
control to use digital I/O instead of a URScript socket (\texttt{92d0ead}) --
the socket approach proved unreliable for real-time gripper control. This
period is also when the VLA approach documented in the older
\texttt{robotic-ai/docs} materials was abandoned in favor of the current
LLM/VLM architecture (Section~\ref{sec:introduction}); no VLA code exists in
this repository's history, meaning the pivot happened before this repository
was created.

Early hardening work in this period: fixing a callback deadlock, workspace
limits, pick-state tracking, plan normalization, a startup-home behavior, and
the real-robot launch file (\texttt{c43c235}); adding threading to the task
interface to avoid a separate deadlock (\texttt{b4817e3},
\texttt{5ab610d}); and introducing horizontal/vertical approach with smooth
linear pick/release trajectories (\texttt{15ddecd}).

\subsection{Session 1: VLM pipeline stand-up and calibration (mid-June 2026)}

Goal: get a real camera into the loop and produce a working
pixel-to-3D-to-robot pipeline.

\textbf{ChArUco detection.} Replaced \texttt{estimatePoseCharucoBoard} with
\texttt{solvePnP} (\texttt{SOLVEPNP\_ITERATIVE}) for more stable pose
estimates; migrated to OpenCV's new \texttt{ArucoDetector} API with
\texttt{CORNER\_REFINE\_NONE}; added manual axis drawing via
\texttt{projectPoints} to support an explicit \texttt{z\_sign} convention.

\textbf{OpenCV 4.5.4 compatibility.} \texttt{cv2.line}/\texttt{circle}/
\texttt{putText} reject \texttt{numpy.int32} scalars on this OpenCV version;
fixed by converting through \texttt{np.round(...).astype(np.int32).tolist()}
everywhere pixel coordinates were drawn.

\textbf{QoS mismatch.} The camera driver publishes with \texttt{BEST\_EFFORT}
(sensor-data QoS); all three camera subscriptions in
\texttt{camera\_services\_node} and \texttt{pixel\_overlay\_viewer\_node} were
still on the default \texttt{RELIABLE} profile and silently received nothing.
Fixed by switching to \texttt{qos\_profile\_sensor\_data}.

\textbf{VLM pick/release state tracking.} In VLM mode, object positions come
from the camera rather than a static scene, so
\texttt{apply\_skill\_result('pick')} -- which only set \texttt{robot.holding}
when \texttt{find\_object\_near} matched a known object -- never fired, and
the plan validator then rejected the paired release for ``nothing held''.
Fixed by making the holding-state guardrails mode-aware: full guardrails
stay active in LLM mode; VLM mode skips object-matching entirely.

\textbf{Workspace limits.} The default \texttt{x: [0.15, 0.85]} excluded
objects near the robot base; corrected to match the real table and camera
coverage (an operator-maintained value in \texttt{scene\_mock.json}, not a
one-time fix).

\subsection{Session 2: code review and refactor}

A systematic pass fixing five concrete bugs surfaced by code review, followed
by deduplication and dead-code removal.

Bugs fixed: a call to a nonexistent \texttt{get\_total\_used\_tokens()}
method (corrected to \texttt{get\_total\_usage()}); a misspelt
\texttt{force\_replanning} kwarg silently swallowed by \texttt{**kwargs}
instead of reaching every reasoning method's \texttt{force\_replan} parameter;
an inconsistent \texttt{"end of simulation"} vs.\ \texttt{"end\_of\_simulation"}
JSON key between the ReAct prompt and its parser; a live \texttt{assert} on
LLM output in Tree-of-Thought batch evaluation that was a guaranteed crash on
any model returning an unexpected count (removed entirely, replaced by
always using IID evaluation); and a ChArUco board-geometry mismatch between
\texttt{config.py}'s defaults ($5\times7$, 0.03/0.015\,m) and the actual
launch-file-authoritative board ($3\times4$, 0.062/0.031\,m).

Refactors: a shared \texttt{\_select\_prompts()} helper collapsing an
identical VLM/LLM prompt-selection branch repeated across six reasoning
methods (11 call sites); a new \texttt{agent\_runner.py} with
\texttt{run\_plan\_loop()} removing $\sim$30 duplicated lines between the LLM
and VLM planner nodes; a \texttt{\_wait\_for\_future()} helper replacing two
bare busy-wait loops with a 10-second-timeout version; and removal of a
byte-identical \texttt{\_call\_nebius} method (Nebius is OpenAI-API-compatible
and only differs by \texttt{base\_url}, so it now routes through
\texttt{\_call\_openai}). Several unused classes and a broken \texttt{\_\_main\_\_}
block were deleted, and all node log prefixes were standardized to
\texttt{[ClassName]}.

\subsection{Session 3: web GUI, emergency stop, debug recorder}

The largest single addition: the full \texttt{robo\_reason\_gui} package,
built across six phases (scaffold $\to$ read-only dashboard $\to$ chat wired
to plan/execute $\to$ live streamed execution $\to$ stack orchestration with
live planner params $\to$ GUI-owned UR driver with auto-retry $\to$ polish),
described in Section~\ref{sec:architecture}.

Alongside the GUI itself:
\begin{itemize}[nosep]
  \item A new \texttt{CancelExecution.srv} service implementing an emergency
    stop: abort the current plan, cancel the active
    \texttt{/execute\_skill} goal, and send the arm home. Wired through
    \texttt{plan\_manager\_node}'s step loop (a \texttt{threading.Event}
    checked before every step), the bridge node (deliberately not gated
    behind the same lock as normal command execution, so it stays callable
    while a plan is blocked mid-execution), and a red Stop button in the
    frontend.
  \item A shared \texttt{\_strip\_json\_fence()} helper across all six
    reasoning methods to tolerate markdown code-fence-wrapped JSON, plus a
    fix for a \texttt{force\_json}/\texttt{force\_json\_response} kwarg-name
    mismatch that had silently disabled JSON enforcement on some call sites.
  \item A model registry refresh: Groq's lineup fully replaced (old Llama
    4/3.3/3.1 models removed; \texttt{qwen3-32b} and vision-capable
    \texttt{qwen3.6-27b} added), Nebius additions, and every launch file's
    hardcoded default model name updated to match (several had gone stale,
    pointing at models no longer in the registry).
  \item A per-run debug recorder (\texttt{debug\_recorder.py}) capturing
    command, config, raw response, logs, and camera frame/overlay per
    \texttt{/plan\_task} call, plus a \texttt{DEBUG\_TIMEZONE} setting so
    debug-folder timestamps match the operator's local time rather than the
    container's UTC clock.
  \item Two failure modes newly surfaced by that debug recorder: Groq empty
    completions on \texttt{fhp}/\texttt{ffhp} (likely exhausting the
    reasoning-token budget on the longer chained prompt), and Nebius
    bad-grounding-depth failures where the plan parses fine but the returned
    pixel lands off the target's surface. A reduced-schema mitigation was
    discussed but not implemented in this session.
\end{itemize}

\subsection{Session 4: VLM$\to$LLM hybrid, grasp-width geometry, grounding fixes}

\textbf{VLM$\to$LLM hybrid pipeline} (\texttt{mode=VLM\_LLM}, new
\texttt{vlm\_llm\_planner\_node.py} + \texttt{scene\_grounder.py}): a VLM
grounding call detects objects/targets as pixel centers, deprojects them to
world coordinates, assembles a generated scene in memory, then hands off to
the standard LLM planner. Grounding and planning model/temperature are wired
independently through config, launch files, and the GUI.

This session also fixed a project-wide \textbf{VLM pixel-coordinate
convention bug}: Qwen-family VLMs ignore a requested \texttt{[h, w]}
(row, col) output order and always emit their native \texttt{[x, y]}
(col, row) convention, causing a systematic axis-swap between the detected
scene and the real one. Every reasoning-method prompt, \texttt{skills.py},
\texttt{extraction\_classes.py}, and both VLM planner nodes were standardized
to consistently use \texttt{[x, y]}.

\textbf{Width-aware TCP offset and release-height fix.} Added
\texttt{grasp\_width} to the action schema and every prompt; added a
piecewise-linear \texttt{tcp\_offset\_z\_for\_width()} calibration table
(since the RG2's fingers pivot, flange-to-contact distance varies with
object width); capped pick descent depth so tall objects are gripped nearer
their top instead of driving the rigid TCP clamp into them; and fixed a
\textbf{double-counted release height} in the hybrid pipeline where a
depth-measured top surface was having the VLM's separate, non-depth-grounded
size estimate added on top of it a second time. See
\texttt{docs/GRASP\_GEOMETRY\_PIPELINE.md} for the full geometry derivation.

\textbf{Further hardening} in the same session: local-plane-centroid
correction for VLM-click grounding error near object edges; a
\texttt{table\_surface\_z}/\texttt{top\_z} sign inversion fix in object-height
computation; moving TCP calibration out of \texttt{config.py} into a
dedicated \texttt{grasp\_geometry.py}; a new terminal-log capture service so
each debug run also captures camera/robot/stack subprocess logs; the UR
driver supervisor fast-failing on an \texttt{[ERROR]} log marker instead of
always waiting a full 30-second timeout; and consolidating scattered
per-class settings into \texttt{config.py} as the single settings source.

A first \textbf{bbox grounding mode} (point vs.\ bounding box for VLM pixel
grounding) and a \textbf{CoT-SC timeout fix} closed out the session: no
provider client had ever had an explicit HTTP timeout (SDK defaults up to
600\,s exceeded the app's 300\,s plan timeout), and \texttt{cot\_sc} ran its
$k{=}5$ independent samples sequentially, multiplying exposure to any single
slow call; fixed with a bounded, configurable
\texttt{Settings.REQUEST\_TIMEOUT\_S} on every client and parallelized
sampling via \texttt{ThreadPoolExecutor} (with a lock added around metrics
updates, now hit concurrently).

Three further robustness fixes closed the session: a \texttt{cot\_sc}
crash when a sampled action came back as a non-dict, now caught and raised as
a catchable \texttt{ActionParsingError}; an unclosed \texttt{<think>} block
(model exhausts its token budget mid-thought) that previously bypassed the
blank-response retry path, now truncated and retried correctly; and a new
\texttt{VLM\_REASONING\_EFFORT} setting forwarding Groq's
\texttt{reasoning\_effort} parameter end-to-end, as a mitigation for
Groq/Qwen's imprecise in-\texttt{<think>} percentage-estimation grounding
(not yet A/B tested on hardware). The full reasoning test suite stood at
45/45 passing at the end of this session.

\subsection{Session 5: release-geometry determinism, zone-collision spacing, and the LLM-vs-VLM benchmark}

The current session, on branch \texttt{feature/refactoring}.

\textbf{Zone-release collision spacing} (\texttt{a1b744a}, committed):
requesting multiple objects into the same zone reliably produced overlapping
releases at the zone's exact center; a new
\texttt{distribute\_zone\_releases()} spaces them along a line-then-grid
pattern, detecting collisions by comparing release positions to each other
\emph{within the plan} (mode-agnostic) rather than against a scene's static
target registry (which only ever matched LLM's hand-authored positions).

\textbf{LLM-mode release geometry made deterministic}: two hardware-reported
bugs traced to the same root cause -- every LLM prompt asked the model to
compute release height/position via its own arithmetic, and that arithmetic
was repeatedly wrong on real hardware. Fixed by computing both
deterministically in code (\texttt{\_fix\_object\_height},
\texttt{\_fix\_release\_height}), mirroring what VLM mode already did.

\textbf{VLM-mode grasp-width sampling fix}: bbox edge pixels sit exactly on
an object's silhouette boundary, the worst place to get valid depth, and a
single failed edge sample aborted the entire batched deprojection call --
crashing an otherwise-fine plan. Fixed by making width sampling a
best-effort, exception-isolated call sampling slightly inset from the raw
edge.

\textbf{Benchmark infrastructure and study}: a full LLM-vs-VLM comparison,
covered in detail in Section~\ref{sec:benchmark}.

Full details of every fix in this section, including exact file lists and
line-level rationale, are in \texttt{docs/SESSION\_CONTEXT.md}.
